\chapter{Emergency Report Form and Contact Numbers}\label{ch:emergency_forms}

\section{Emergency Report Form}

Copy this form, fill in the family information, and keep it with your first aid kit. During an emergency it helps you record what matters and hand over a complete report to responders.

\textbf{Personal Information}

\begin{description}
  \item[Name:] \rule{0.6\textwidth}{0.4pt}
  \item[Date of birth:] \rule{0.55\textwidth}{0.4pt}
  \item[Address:] \rule{0.58\textwidth}{0.4pt}
  \item[Phone:] \rule{0.6\textwidth}{0.4pt}
  \item[Emergency contact:] \rule{0.42\textwidth}{0.4pt}
  \item[Family doctor:] \rule{0.5\textwidth}{0.4pt}
  \item[Insurance / health number:] \rule{0.35\textwidth}{0.4pt}
\end{description}

\textbf{Medical Information}

\begin{description}
  \item[Allergies:] \rule{0.7\textwidth}{0.4pt}
  \item[Current medications:] \rule{0.52\textwidth}{0.4pt}
  \item[Medical conditions:] \rule{0.55\textwidth}{0.4pt}
  \item[Last tetanus booster:] \rule{0.48\textwidth}{0.4pt}
  \item[Blood type (if known):] \rule{0.45\textwidth}{0.4pt}
\end{description}

\textbf{Incident Report (complete at the scene)}

\begin{description}
  \item[Date and time of incident:] \rule{0.42\textwidth}{0.4pt}
  \item[Time emergency services called:] \rule{0.32\textwidth}{0.4pt}
  \item[What happened:] \rule{0.68\textwidth}{0.4pt}
  \item[Location / address:] \rule{0.55\textwidth}{0.4pt}
  \item[How the person was found:] \rule{0.5\textwidth}{0.4pt}
  \item[Level of consciousness (AVPU):] \rule{0.45\textwidth}{0.4pt}
  \item[Breathing (rate and quality):] \rule{0.45\textwidth}{0.4pt}
  \item[Pulse (rate and rhythm, if checked):] \rule{0.38\textwidth}{0.4pt}
  \item[Observations (SAMPLE/OPQRST):] \rule{0.4\textwidth}{0.4pt}
  \item[Time of last food/drink:] \rule{0.45\textwidth}{0.4pt}
\end{description}

\textbf{Interventions Provided}

\begin{description}
  \item[First aid given:] \rule{0.65\textwidth}{0.4pt}
  \item[Medications given (dose and time):] \rule{0.4\textwidth}{0.4pt}
  \item[Time CPR started:] \rule{0.5\textwidth}{0.4pt}
  \item[Time AED attached / shocks delivered:] \rule{0.3\textwidth}{0.4pt}
  \item[Other notes:] \rule{0.72\textwidth}{0.4pt}
\end{description}

\quicknote{\textbf{How to report to responders (SBAR):} Situation -- what is happening now. Background -- what happened. Assessment -- what you found. Recommendation -- what was done and what is needed. See Section~\ref{sec:sbar} for details.}

\section{International Emergency Numbers}

Emergency numbers vary by country. Program the local number into your phone when you travel.

\begin{center}
\begin{tabularx}{\textwidth}{@{}>{\raggedright\arraybackslash}l>{\raggedright\arraybackslash}l>{\raggedright\arraybackslash}X@{}}
\hline
\textbf{Country} & \textbf{Emergency Number} & \textbf{Notes} \\
\hline
United States & 911 & Police, fire, medical \\
Canada & 911 & Police, fire, medical \\
United Kingdom & 999 or 112 & Police, fire, medical \\
European Union & 112 & Single European emergency number \\
Australia & 000 & Police, fire, medical \\
New Zealand & 111 & Police, fire, medical \\
India & 112 & National unified emergency number \\
India & 108 & Emergency ambulance (state dial-access services) \\
India & 102 & Ambulance (in many states) \\
India & 100 & Police \\
India & 101 & Fire \\
India & 104 & Health helpline (in many states) \\
India & 1078 & National Disaster Management Authority (NDMA) \\
China & 120 & Medical; 110 police; 119 fire \\
Japan & 119 & Fire and medical; 110 police \\
Singapore & 995 & Medical \\
South Korea & 119 & Fire and medical \\
Mexico & 911 & Police, fire, medical \\
South Africa & 10177 & Medical \\
\hline
\end{tabularx}
\end{center}

\textbf{Poison Control}: In India, the National Poisons Information Centre (NPIC), All India Institute of Medical Sciences (AIIMS), New Delhi, operates a 24-hour helpline at \textbf{011-26589391} (also \textbf{011-26593677}). United States 1-800-222-1222 (24/7). In most countries, a poison control center can be reached through the national emergency number or a dedicated hotline -- store the number in your phone and on the family form.

\textbf{WARNING:} \textbf{When you call emergency services, be ready to state:} your exact location, what happened, how many people are involved, the condition of the person (responsive? breathing?), and what is being done. Do not hang up until the dispatcher tells you to.

\section{Medical Alert Identification}

For a person with severe allergies, serious medical conditions, or important medications, a medical alert ID (bracelet, necklace, or card) can save critical time:

\begin{itemize}[leftmargin=*]
  \item Check the wrist, neck, and ankle of an unresponsive person for a medical alert ID
  \item The ID should list the condition, allergies, and an emergency contact number
  \item Tell responders exactly what the ID says -- do not remove it unless asked
  \item Anyone who is unresponsive, or unable to speak for themselves, should consider wearing one
\end{itemize}

\section{Review Questions}

\begin{enumerate}[leftmargin=*]
\item Why is it important to record the exact time emergency services were called?
\item Which four pieces of information must you be ready to give when calling emergency services?
\item Why should a medical alert ID never be removed by a first aider?
\end{enumerate}
\index{emergency report form}
\index{emergency numbers}
\index{medical alert ID}
\index{poison control}
